% Exact proof additions for C1.04, independently audited 2026-11-19.

% Replace the completion paragraph by this clarification or integrate it.
In the classifying-space construction one takes the classifying space of the \emph{entire groupoid of chains} in each degree, with isomorphisms at every vertex and the commuting identifications on each arrow. Replacing only the object set by classifying spaces while leaving the original arrow set unchanged would not suffice. For example, for a single object with automorphism group \(G\), the groupoid of length-\(n\) strings is
\[
 G^n\mathbin{/\mkern-6mu/}G^{n+1}.
\]
Its action changes the source and target trivializations at every vertex; it is transitive with stabilizer the diagonal copy of \(G\). Its classifying space is therefore \(BG\) in every degree, and the degeneracy from degree zero to the equivalence subspace in degree one is an equivalence. This is the Rezk classifying-diagram description. Applying the same construction to the full incidence diagram, with homotopy fibre products for compositions, gives the first complete-Segal model; Rezk completion of the raw action-category nerve gives the same model.

% Add before invoking chartwise equivariance.
Choose the compatible boundary fibrations and collar data once on the orbifold resolution, and pull them back to each atlas chart. They can be constructed equivariantly: the normal-crossings blow-ups preserve the incidence labels, invariant metrics and connections exist by finite-group averaging and orbifold partitions of unity, and the iterated collar construction is performed successively on the preserved boundary fibrations. Thus a change of chart preserves the pulled-back structure. One is not asserting that independent arbitrary collar choices in different charts agree. The chartwise source maps are representable proper locally trivial bundles. Their pullbacks therefore model homotopy fibre products; after passing to an atlas and its chain groupoids, the Segal fibre-product comparisons remain equivalences of homotopy types.

% Add to the definition of orbifold Exit by atlas descent.
This atlas definition is independent of the atlas. To see the necessary descent statement first for a surjective stratum-preserving local diffeomorphism \(p:V\to U\), cover \(U\) by open sets on which \(p\) has a stratum-preserving local section. On any such open set the augmented \v Cech diagram of \(p\) has an extra degeneracy supplied by the section, so its image under the exit functor is a split augmented simplicial diagram. Its colimit is the exit category of that open set. Exit open-cover excision, \cite[Theorem A.7.1]{LurieHA}, and interchange of colimits globalize this equivalence to \(U\). Given two orbifold atlases, their fibre product is a common refinement; the preceding assertion applied in each direction of its bisimplicial atlas shows that both colimits are equivalent. On a quotient chart \([U/G]\), the atlas nerve is \(U\times G^\bullet\), so its exit colimit is precisely the homotopy-orbit category \(\operatorname{Exit}(U)\mathbin{/\mkern-6mu/}G\), including ineffective isotropy.
